\documentclass[11pt,a4paper]{article}
\usepackage[UTF8,heading=true]{ctex}
\usepackage[a4paper,left=25mm,right=25mm,top=22mm,bottom=22mm]{geometry}
\usepackage{fontspec,xeCJK,amsmath,amssymb,booktabs,array}
\usepackage[hidelinks]{hyperref}
\setmainfont{Times New Roman}
\setCJKmainfont{SimSun}[BoldFont=SimHei]
\setCJKsansfont{SimHei}
\setlength{\parindent}{2em}
\setlength{\parskip}{3pt}
\linespread{1.18}
\emergencystretch=1.5em
\widowpenalty=10000
\clubpenalty=10000
\ctexset{section={format=\Large\bfseries},subsection={format=\large\bfseries}}
\pagestyle{plain}

\begin{document}

\section*{摘\quad 要}

本文研究含光伏、储能和外部电网的小区微网购电问题。难点在于光伏与负载随时间变化，电池存在容量、功率和充放电损耗，而普通购电、改单和紧急购电又采用不同的价格规则。为使四问前后连贯，本文先建立统一的十分钟电量流和储能递推关系，再根据各问的信息条件逐步加入供需预测、风险准备、日内预报和波动电价。

针对问题一，建立\textbf{全时段电量流线性规划模型}，统筹光伏直供、外网购电和储能移峰，并用日初、日末储电量相等保证费用比较公平。所得代表日费用为35126.85元，低于峰谷购电策略的38276.64元和光伏空间预留策略的37944.35元。

针对问题二，建立\textbf{周期预测--近期加权准备--跨日线性规划模型}。负载与光伏分别依据历史规律预测，再用近期净需求误差确定安全准备量，以应对五倍价格的紧急购电。2025年2月至12月总费用为13583810.84元，其中普通电费12979560.57元、紧急电费604250.27元；相较等权准备方案节省44.66万元。参数检验表明，0.7的日期衰减在快速适应近期变化和保留足够历史信息之间取得了较好平衡。

针对问题三，建立\textbf{历史融合预报--动态安全准备--日内滚动规划模型}。模型将官方光伏预报与近期历史按过往误差自动融合，并在0:00、6:00、12:00和18:00滚动调整尚未执行的订单。全年总费用为13146535.49元，比不加准备量的初始融合方案节省94.22万元。在日内反馈和改单权限相同的条件下，引入6:00、12:00和18:00的新预报又节省21.78万元，说明后发预报值得使用，但每次是否改单仍需结合调整费和储电状态判断。

针对问题四，将未知波动电价纳入预测和调度。零点一次订电分支采用\textbf{滚动提升树--零点线性规划模型}，总费用为14248570.87元；日内调整分支采用\textbf{滚动岭回归--日内调整线性规划模型}，总费用为13822901.35元。后者虽未取得候选中的绝对最低费用，却以0.016\%的费用差换得更少的紧急购电和更清楚的价格修正依据。逐段回算表明，各方案满足供需平衡、储能容量和充放电功率限制。模型能够在费用、供电风险和解释性之间形成清楚的取舍，可为社区、校园和园区微网的滚动购电提供参考。

\noindent\textbf{关键词：}微网调度；储能优化；滚动规划；风险准备；波动电价

\newpage
\section{问题分析}

\subsection{问题一：已知一天曲线时，电量应当怎样跨时段转移}

问题一给出了完整的日内电价、负载和光伏预测，困难并不在于猜测未来，而在于安排电量何时买、何时存、何时用。光伏最多的时段未必是负载最高的时段，最低电价时段也未必有足够储能空间。若只在低价时尽量充电，电池可能在午间光伏到来前已经装满；若过早放电，又可能在晚间高价和高负载同时出现时失去缓冲能力。因此，需要把全天144个十分钟时段放在一起优化，而不能逐段作出彼此独立的选择。

日初与日末储电量相等是本问的重要条件。它阻止模型通过消耗初始库存制造“今天少买电”的表面低费用，使各方案在相同起点和终点下比较。基于这一结构，问题一采用\textbf{全时段电量流线性规划模型}，将外网直供、外网充电、光伏充电、电池放电和储电量变化统一写入一个优化问题。

\subsection{问题二：计划在零点确定，真实供需随后才出现}

问题二的普通购电必须在每天零点确定，而当天真实负载和光伏还没有发生。附件中的全年数据只能作为事后运行记录和此前历史，不能当作零点已经知道的未来。负载有较明显的星期规律，光伏更受近期天气与季节变化影响，因此二者需要分开预测；预测完成后再合成为净需求，才能保留各自的现实含义。

本问的费用具有明显的不对称性。多订的普通电即使用不完也要付款，少订后真正缺电的部分却要按五倍电价紧急购买。电池能吸收一部分偏差，但会受到最低储电量和放电功率限制，因而“全天总电量够”并不代表每个时段都不会缺电。模型既要给出通常情况下的供需判断，又要依据近期发生过的低估误差留出适量准备。由此采用\textbf{周期预测--近期加权准备--跨日线性规划模型}，并让日末库存自然传给下一天，使当前订电不会把风险简单推迟到次日。

\subsection{问题三：新预报是否值得用，取决于它能否覆盖改单成本}

问题三允许在6:00、12:00和18:00得到新的光伏预报后调整未来订单。新增信息有助于纠正零点判断，但修改订单需要支付调整费，电池也可能已经足以消化预测变化。因此，“获得新预报”和“必须改单”是两件事。合理策略应当先更新未来净需求和实际库存，再比较维持原订单、调增、调减和可能紧急购电的代价。

附件3按小时给出光伏预报，而购电结果按十分钟填写。相邻整点之间采用线性插值，使预报与144个运行时段对应；附件2中的实际光伏只在时段发生后用于更新误差和库存。官方预报反映当天信息，近期同刻历史提供稳定参照，两者各有作用。因而采用\textbf{历史融合预报--动态安全准备--日内滚动规划模型}，并在每个发布时间只调整尚未执行的区间。评价后发预报时，还要把“允许日内反馈和改单”的收益与“新增官方预报”的收益分开，否则会夸大预报本身的贡献。

\subsection{问题四：未来价格也未知后，购电时点本身需要重新判断}

问题四把固定电价改为波动电价。电价变化不只是把原订单重新计价，还会改变什么时候适合充电、什么时候应保留电量以及什么时候放电。因此，需要先用决策时已经发生的价格和历史规律判断未来价格，再把预测价格送入购电与储能规划。

两条分支的权限不同。问题四-2仍在零点一次确定当天普通订单，无法利用当天已发生价格修正，因此更需要表达历史价格、净负荷、星期和季节之间的非线性组合，采用\textbf{滚动提升树--零点线性规划模型}。问题四-3可以在日内更新，价格特征之间相关性较强，又需要解释每次修正来自哪里，因此采用\textbf{滚动岭回归--日内调整线性规划模型}。两种价格模型都只提供决策时的未来判断，最终费用仍按实际电价和真实运行结果结算。

\section{模型评价与推广}

\subsection{模型优点与创新}

\textbf{（1）四问共用一套可以核对的电量流骨架。}本文始终用十分钟电量连接负载、光伏、外网和电池，并统一处理90\%充电效率、90\%放电效率、储电上下限及功率限制。各问只在信息条件和费用规则发生变化时增加相应模块，避免同一物理过程被重复解释成不同含义。问题一的确定性调度由线性规划完成\cite{dantzig}，后续各问沿用同一基础关系，便于逐段复算。

\textbf{（2）把预测和高价缺口风险分开处理。}负载与光伏分别预测，使星期用电规律和近期光伏变化都能被解释\cite{hyndman}；误差准备则专门处理五倍紧急电价造成的损失不对称。候选比较显示，提升树即使改善部分预测误差，也没有超过带近期准备的完整方案。该设计说明模型追求的是最终运行费用和供电可靠性，而不是单一的拟合指标。

\textbf{（3）把日内信息的价值拆成两个可说明的部分。}问题三没有把197.28万元的总差额全部归功于新预报，而是先计算日内反馈与改单节省的175.51万元，再计算后发官方预报额外节省的21.78万元。这样可以直接回答题目关于日内预报是否值得引入的问题，也避免把权限变化误写成预测方法的贡献。

\textbf{（4）根据决策方式选择价格模型。}问题四没有机械地给两条分支使用同一种算法。零点分支采用提升树表达非线性关系\cite{friedman}；日内分支采用岭回归稳定处理高度相关的历史价格特征\cite{hoerl}。后者以极小费用差换取更少紧急购电和更清楚的系数含义，体现了经济性、风险和可解释性的综合选择。

\textbf{（5）结论能够追溯到实际运行过程。}本文使用SciPy和scikit-learn完成数值计算与候选预测\cite{virtanen,pedregosa}，线性规划由HiGHS求解\cite{huangfu}。全年费用均由保存的普通订单、紧急购电和储电轨迹重新计算，并逐段检查供需、库存、功率和结算关系。指定日期的图表进一步说明电池何时补缺、何时达到下限以及为什么仍会发生紧急购电，使模型结论既有总量证据，也有运行过程支撑。

\subsection{模型局限与改进方向}

首先，历史窗口、误差分位、日期衰减和价格模型参数主要在2025年数据内比较。现有结果能够说明本题数据范围内的取舍，但不能证明同一组参数在其他年份或其他地区仍然最优。若获得跨年数据，可按年份滚动回测，并在不同季节分别校准参数，以检验真正的跨年泛化能力。

其次，本文先预测负载、光伏和价格，再把预测结果送入优化。这样的分层结构容易解释，也便于定位误差来源，但没有同时描述三类不确定量可能共同出现的极端情景。后续可构造少量有现实含义的联合情景，比较连续阴天、高负载和高电价叠加时的计划，进一步控制尾部风险。

再次，题目没有提供电池老化费用、外送电价、网络损耗和可中断负载信息，因此本文未把这些因素写入目标函数。若用于真实微网，可增加循环寿命成本、需求响应和售电约束；末端库存的参考价值也应根据下一日市场规则重新估计，而不能直接视为可出售资产。

\subsection{模型推广性评估}

模型的核心输入只有分时负载、本地可再生能源、设备参数、交易价格和允许的决策时点。替换这些数据并重新校准预测与准备参数后，可用于社区、校园、工业园区及偏远地区的微网运行，也可扩展到多台储能、风光互补、充电站和可调负载。滚动求解方式允许系统随着新数据修正未来计划，适合部署在需要持续更新的能源管理平台中。

从社会应用看，该方法优先消纳本地光伏，并通过储能和合理购电减少高峰期对外网的额外压力，有助于提高可再生能源利用率、降低能源浪费，服务碳达峰和碳中和目标\cite{carbonplan,energyplan}。推广时还应把安全供电和居民基本用电放在成本优化之前，使节费收益建立在可靠、公平的能源服务之上。这样的应用方向既符合绿色发展和资源节约要求，也体现以人民为中心、共建共享的发展取向。

\newpage
\begin{thebibliography}{99}
\bibitem{dantzig} Dantzig G B. Linear Programming and Extensions[M]. Princeton: Princeton University Press, 1963.
\bibitem{hoerl} Hoerl A E, Kennard R W. Ridge regression: Biased estimation for nonorthogonal problems[J]. Technometrics, 1970, 12(1): 55--67.
\bibitem{friedman} Friedman J H. Greedy function approximation: A gradient boosting machine[J]. The Annals of Statistics, 2001, 29(5): 1189--1232.
\bibitem{huangfu} Huangfu Q, Hall J A J. Parallelizing the dual revised simplex method[J]. Mathematical Programming Computation, 2018, 10(1): 119--142.
\bibitem{hyndman} Hyndman R J, Athanasopoulos G. Forecasting: Principles and Practice[M]. 3rd ed. Melbourne: OTexts, 2021.
\bibitem{virtanen} Virtanen P, Gommers R, Oliphant T E, et al. SciPy 1.0: Fundamental algorithms for scientific computing in Python[J]. Nature Methods, 2020, 17: 261--272.
\bibitem{pedregosa} Pedregosa F, Varoquaux G, Gramfort A, et al. Scikit-learn: Machine learning in Python[J]. Journal of Machine Learning Research, 2011, 12: 2825--2830.
\bibitem{carbonplan} 国务院. 2030年前碳达峰行动方案[Z]. 2021.
\bibitem{energyplan} 国家发展改革委, 国家能源局. “十四五”现代能源体系规划[Z]. 2022.
\end{thebibliography}

\end{document}
